% !TEX root = ../glossario.tex
% Lettera B del Glossario (vedi capitoli/a.tex per le istruzioni).
\phantomsection
\letteraglossario{B}

\termine{Backoffice}{Area amministrativa di un'applicazione, riservata a chi gestisce
  contenuti, utenti, configurazioni e processi; non è visibile agli
  utenti finali.}

\termine{Baseline}{Versione consolidata e approvata di un prodotto o di un insieme di artefatti in un determinato istante temporale, che funge da punto di riferimento per lo sviluppo futuro e può essere modificata solo attraverso procedure formali di controllo.}

\termine{Best practice}{Procedura o metodo riconosciuto come il modo più efficiente ed efficace per svolgere un compito o affrontare un problema.}

\termine{Branch}{Linea di sviluppo parallela all'interno di un sistema di controllo di versione (come Git), che permette di lavorare su nuove funzionalità o correzioni in modo isolato senza intaccare il ramo principale (main).}

\termine{Bucket}{Nel capitolato C2, contenitore configurabile di fonti documentali, con
  nome, descrizione e destinatari autorizzati: stabilisce chi può
  consultare la conoscenza ricavata da quelle fonti.}

\termine{Bug}{Errore, anomalia o difetto nel codice sorgente di un software che ne causa il malfunzionamento o produce un risultato inatteso.}

\termine{Business Intelligence}{Insieme di processi e strumenti che raccolgono, organizzano e
  analizzano i dati di un'azienda per presentarli in forma utile alle
  decisioni, per esempio con report, indicatori e dashboard.}
